\section{HUM duality theory}
\label{sec:hum-duality}

\subsection{The finite-dimensional (ODE) prototype}
\label{ssec:ode-prototype}

The infinite-dimensional theory below generalizes a finite-dimensional
template for the linear ODE system
\begin{equation}
  \label{eq:ode-system}
  y'(t) = Ay(t) + Bv(t) \ \text{in } (0,T), \qquad y(0) = y_0,
\end{equation}
with $n,m \in \mathbb{N}$ (the state and control dimensions),
$A \in \R^{n\times n}$, $B \in \R^{n\times m}$, $y : [0,T]\to\R^n$,
$v \in L^2(0,T;\R^m)$. This template is stated here in full because it is the
direct prerequisite for the discrete/ODE uniform-controllability theory of
Section~\ref{sec:discrete}: a space semi-discretization reduces the PDE
control problem to exactly this ODE setting.

\begin{theorem}[Kalman rank condition]
  \label{thm:kalman}
  Let $K = [B, AB, \dots, A^{n-1}B]$. System \eqref{eq:ode-system} is exactly
  controllable in every time $T>0$ if and only if $\rank K = n$.
\end{theorem}
\begin{proof}[Source]
  \citet[Thm.~1.16]{coron2007control}.
\end{proof}

The adjoint system of \eqref{eq:ode-system} is
\begin{equation}
  \label{eq:ode-adjoint}
  \varphi'(t) = -A^*\varphi(t) \ \text{in } (0,T), \qquad \varphi(T) = \varphi_T.
\end{equation}

\begin{theorem}[Controllability $\Leftrightarrow$ observability]
  \label{thm:ode-observability}
  System \eqref{eq:ode-system} is controllable in time $T$ if and only if the
  adjoint system \eqref{eq:ode-adjoint} is \emph{observable} in time $T$, i.e.\
  there exists $C=C(T)>0$ such that, for every solution $\varphi$ of
  \eqref{eq:ode-adjoint},
  \begin{equation}
    \label{eq:ode-observability-ineq}
    |\varphi(0)|^2_{\R^n} \le C\int_0^T |B^*\varphi(t)|^2_{\R^m}\,dt.
  \end{equation}
  Both properties hold in every time $T$ if and only if the Kalman rank
  condition (Theorem~\ref{thm:kalman}) holds.
\end{theorem}
\begin{proof}[Source]
  \citet[Thm.~2.1.1]{zuazua2002controllability}.
\end{proof}

\begin{theorem}[HUM control, finite-dimensional]
  \label{thm:ode-hum}
  Let $\mathcal{J}: \R^n \to \R$,
  $\mathcal{J}(\varphi_T) = \tfrac12\int_0^T |B^*\varphi|^2 + (y_0,\varphi(0))$,
  where $\varphi$ solves \eqref{eq:ode-adjoint} with $\varphi(T)=\varphi_T$.
  If $\mathcal{J}$ has a minimizer $\hat\varphi_T$, then, letting $\hat\varphi$
  solve \eqref{eq:ode-adjoint} with $\hat\varphi(T)=\hat\varphi_T$, the control
  $v = B^*\hat\varphi$ is the minimal-$L^2$-norm control steering
  \eqref{eq:ode-system} from $y_0$ to $0$ at time $T$.
\end{theorem}

The construction hinges on the \emph{Gramian operator}
\begin{equation}
  \label{eq:ode-gramian}
  \Lambda: \R^n \to \R^n, \qquad
  \Lambda(g) := \int_0^T B^*\varphi_g(s)\,ds,
\end{equation}
where $\varphi_g$ solves \eqref{eq:ode-adjoint} with final datum $g$; $\Lambda$
is well defined by duality \citep{boyer2013penalised}. This is the same
Gramian apparatus used below for the PDE-level construction, over $\R^n$
instead of an infinite-dimensional Hilbert space.

\subsection{Infinite-dimensional generalization}
\label{ssec:hum-infinite}

In the PDE setting the Kalman rank condition has no direct analogue; more
delicate tools are needed, and in general exact controllability of the heat
equation fails while approximate controllability holds
\citep{zuazua2002controllability} -- see also \citet{coron2007control,
zabczyk2020mathematical} for the general control-theoretic background this
document draws on. Approximate controllability of the 1D heat equation can
also be established directly via a variational approach
\citep{fabre1995approximate}, as an alternative to the HUM route developed
below. Nonetheless the HUM construction extends: with $H$, $U$ the state and
control spaces of Section~\ref{sec:wellposedness}, and adjoint system in $H$,
\begin{equation}
  \label{eq:pde-adjoint}
  \varphi'(t) = -A^*\varphi(t) \ \text{in } (0,T), \qquad \varphi(T) = g \in H,
\end{equation}

\subsubsection*{Internal control}

\begin{proposition}
  \label{prop:pde-gramian}
  The Gramian $\Lambda: H \to L^2(0,T;U)$ is a positive-definite, symmetric,
  bounded linear operator \citep{boyer2013penalised}, giving rise to the
  bilinear form and dual functional
  \begin{equation}
    \label{eq:pde-functional}
    a(g,\xi) = \int_0^T B^*\varphi_g \cdot B^*\varphi_\xi\,dt, \qquad
    J(g) = \tfrac12 a(g,g) + \langle y_0, \varphi_g(0)\rangle,
  \end{equation}
  with $\langle\cdot,\cdot\rangle$ the $H\times H'$ duality pairing. For
  internal control, $H = U = L^2(0,1)$ and $B^* = B = \mathbf{1}_\omega$.
\end{proposition}

The thesis develops \eqref{eq:pde-functional} directly and moves to
\emph{penalized} HUM (Section~\ref{sec:penalized}) as its route to a
well-posed minimization, rather than stating a standalone
existence/uniqueness theorem for the unpenalized functional.

\begin{remark}[Gap in the thesis, now filled independently -- not from the
  thesis]
  \label{rem:hum-existence}
  The classical existence/uniqueness argument for the unpenalized
  infinite-dimensional minimizer is the original Lions completion-space
  construction. Let
  $F := $ the completion of $H$ with respect to the seminorm
  $\|g\|_F := a(g,g)^{1/2}$. If this seminorm is genuinely a norm on $H$ --
  equivalent to a unique-continuation property for \eqref{eq:pde-adjoint}
  ($B^*\varphi_g \equiv 0$ on $(0,T) \Rightarrow g=0$) -- then $H$ embeds
  continuously and densely into $F$, the linear form
  $g \mapsto -\langle y_0,\varphi_g(0)\rangle$ extends continuously to $F$,
  and Riesz representation on the Hilbert space $F$ yields a unique
  $\hat g \in F$ minimizing $J$ over $F$. This formalizes what the thesis's
  Remark~\ref{rem:coercivity-space} (Section~\ref{sec:known-issues}) describes
  informally as ``the space where $J$ is coercive.''
  Citation: \citet{lions1988controlabilite} -- the original HUM source. This
  is the first citation of Lions' book \emph{directly} in this document; the
  thesis itself cites it only indirectly, via \citealp{russell1990j}'s review
  of it. As with Remarks~\ref{rem:energy-wellposedness} and~\ref{rem:smoothing},
  no specific page/theorem number is claimed here.
\end{remark}

\begin{theorem}[Observability inequality, internal control]
  \label{thm:internal-observability}
  If there exists $C>0$ such that
  \begin{equation}
    \label{eq:internal-observability}
    |\varphi(\cdot,0)|^2_{L^2(0,1)} \le
    C \int_0^1\!\int_0^T |\mathbf{1}_\omega\varphi|^2\,dx\,dt,
  \end{equation}
  then $J$ is coercive.
\end{theorem}

\begin{remark}
  The proof of \eqref{eq:internal-observability} itself requires Carleman
  estimates and is deferred to \citet{fursikov1996imanuvilov}.
\end{remark}

\begin{remark}[Gap in the thesis, now filled independently -- not from the
  thesis: an alternative route]
  \label{rem:lebeau-robbiano}
  Carleman estimates are not the only route to
  \eqref{eq:internal-observability}. The Lebeau--Robbiano spectral inequality
  \[
    \|u\|^2_{L^2(\Omega)} \le C e^{C\lambda}\|u\|^2_{L^2(\omega)}
  \]
  -- valid for any finite combination $u$ of Dirichlet-Laplacian
  eigenfunctions with eigenvalues $\le\lambda$, and any open $\omega\subset
  \Omega$ -- yields null controllability of the heat equation from any open
  subset, in any time $T>0$, via an iterative/telescoping argument
  \citep{lebeau1995controle}. This citation, like
  Remarks~\ref{rem:energy-wellposedness}, \ref{rem:smoothing}
  and~\ref{rem:hum-existence}, is standard-form and was not verified
  page-by-page against the original paper.
\end{remark}

\subsubsection*{Boundary control}

Per Remark~\ref{rem:boundary-h10}, the boundary dual functional
$J_b(\varphi_T)$ lives over $\varphi_T \in H^1_0(0,1)$ rather than
$L^2(0,1)$. No further boundary-specific infinite-dimensional HUM theorem is
stated beyond the general development above.
